\chapter{Anal Fistula Surgery}

\section*{What Is This Surgery?}
Anal fistula surgery is a procedure performed to treat an abnormal, epithelialized tract that connects the anal canal or rectum to the perianal skin. This condition, known as a fistula-in-ano, typically arises from an infection of the anal glands located at the dentate line, leading to an anorectal abscess that subsequently drains and fails to heal completely. The surgical intervention aims to eradicate the fistula tract, eliminate associated sepsis, and prevent recurrence, all while meticulously preserving the integrity of the anal sphincter muscles to maintain fecal continence. The choice of surgical technique is highly individualized, depending on the fistula's complexity, its anatomical relationship to the sphincter, and the patient's overall health.

\section*{Alternative Names}
\begin{itemize}
  \item Fistulotomy
  \item Fistulectomy
  \item Seton placement or seton drainage
  \item Anorectal fistula repair
  \item Endorectal advancement flap for fistula
  \item Ligation of intersphincteric fistula tract (LIFT procedure)
  \item Anal fistula plug insertion
\end{itemize}

\section*{Relevant Surgical Anatomy}
A thorough understanding of anorectal anatomy is paramount for successful anal fistula surgery. The anal canal, approximately 3-4 cm in length, extends from the anorectal ring to the anal verge. It is surrounded by a complex of sphincter muscles crucial for continence. The **internal anal sphincter** is an involuntary smooth muscle, a continuation of the circular muscle of the rectum. The **external anal sphincter** is a voluntary striated muscle, composed of three parts: subcutaneous, superficial, and deep, which fuse with the puborectalis muscle. The **puborectalis muscle**, part of the levator ani complex, forms a sling around the anorectal junction, maintaining the anorectal angle.

The **dentate line**, located roughly midway in the anal canal, marks the transition from columnar rectal epithelium to modified squamous epithelium. At this line, 8-10 anal crypts (crypts of Morgagni) are found, into which anal glands drain. Infection of these glands is the primary cause of most anal fistulas. The **intersphincteric groove** is a palpable depression between the internal and external sphincters, often indicating the level of the dentate line and a common site for the internal opening of a fistula.

Fistulas are classified based on their relationship to the sphincter complex:
\begin{itemize}
  \item \textbf{Intersphincteric:} Tract passes between the internal and external sphincters.
  \item \textbf{Trans-sphincteric:} Tract passes through both internal and external sphincters.
  \item \textbf{Suprasphincteric:} Tract originates in the intersphincteric plane, courses superiorly over the puborectalis, and then descends through the ischioanal fossa to the skin.
  \item \textbf{Extrasphincteric:} Tract originates from the rectum or perianal skin, bypassing the sphincter complex entirely, often due to trauma, Crohn's disease, or malignancy.
\end{itemize}
The **ischioanal (ischio-rectal) fossae** are fat-filled spaces lateral to the anal canal, frequently involved in abscess formation and fistula extension. The **pudendal nerve** provides sensory innervation to the perianal skin and motor innervation to the external anal sphincter, making its preservation critical. The **inferior rectal arteries and veins** supply and drain the anal canal and perianal region. Lymphatic drainage follows the inferior rectal vessels to the inguinal lymph nodes.

\section*{Surgical Principle \& Rationale}
The fundamental principle of anal fistula surgery is to eradicate the chronic inflammatory tract while preserving the function of the anal sphincter muscles to prevent fecal incontinence. The rationale for surgical intervention stems from the fact that an established, epithelialized fistula tract rarely heals spontaneously and serves as a persistent source of infection, pain, and drainage.

The operative concept involves identifying the internal opening, tracing the entire fistula tract, and then either laying it open (fistulotomy) or excising it (fistulectomy) if it involves minimal sphincter muscle. For fistulas that traverse a significant portion of the sphincter complex, sphincter-preserving techniques are employed. These methods aim to close the internal opening and/or drain the tract without dividing the muscle, thereby minimizing the risk of incontinence. Technical goals include precise identification of the internal opening, complete drainage of any associated abscesses, and ensuring the wound heals from the base outwards to prevent premature skin closure and recurrence. The choice of technique is a careful balance between achieving a high cure rate and safeguarding continence.

\section*{History \& Surgical Evolution}
The treatment of anal fistulas has a rich history, with descriptions dating back to ancient civilizations. Hippocrates, in the 5th century BC, described the use of a seton, a thread passed through the fistula tract, for drainage and gradual division. For centuries, surgical approaches were often crude, involving wide excisions or cauterization, frequently resulting in significant morbidity, including incontinence.

The modern era of anal fistula surgery began in the late 19th and early 20th centuries with a better understanding of anal anatomy. In 1934, Milligan and Morgan described the technique of fistulotomy for low fistulas, emphasizing the importance of laying the tract open. A pivotal advancement came in 1976 with the classification system proposed by Alan Parks, which categorized fistulas based on their relationship to the anal sphincter complex. This classification provided a standardized framework for diagnosis and guided surgical decision-making, particularly in distinguishing between low and high fistulas. The recognition that high fistulas required sphincter-preserving techniques led to the development and refinement of procedures such as the endorectal advancement flap, first described in the 1980s. More recently, biological and regenerative approaches, including fibrin glue, collagen plugs, and the Ligation of Intersphincteric Fistula Tract (LIFT) procedure, have emerged, offering less invasive, sphincter-sparing options, particularly for complex or recurrent cases.

\section{Clinical Indications}
Anal fistula surgery is indicated for symptomatic anal fistulas that fail to resolve with conservative management. Key indications include:
\begin{itemize}
  \item Persistent purulent, serous, or bloody drainage from an external opening near the anus.
  \item Recurrent perianal abscesses or persistent drainage after incision and drainage of an abscess.
  \item Chronic pain, discomfort, itching, or irritation in the perianal region attributable to the fistula.
  \item A palpable fibrous cord or an identifiable internal opening on clinical examination or imaging.
  \item Fistulas associated with inflammatory bowel disease, particularly Crohn's disease, that are symptomatic and not responding to medical therapy alone.
  \item Complex fistulas, such as horseshoe fistulas or those with multiple tracts, requiring surgical intervention for definitive management.
  \item Suspected cryptoglandular fistula that has failed to heal spontaneously or with initial non-surgical measures.
\end{itemize}

\section*{Clinical Positioning}
Anal fistula surgery is generally considered the primary curative treatment for established, symptomatic anal fistulas. Unlike acute perianal abscesses, which may resolve with simple drainage, an epithelialized fistula tract rarely heals spontaneously and requires surgical intervention to eliminate the tract and prevent recurrence. Therefore, for most cryptoglandular fistulas, surgery is the first-line definitive therapy once the acute infection has been controlled.

However, its positioning can be secondary or adjunct in specific scenarios:
\begin{itemize}
  \item \textbf{Crohn's Disease:} In patients with Crohn's disease, medical therapy with antibiotics and biologic agents (e.g., anti-TNF-alpha agents) is often attempted first to control underlying inflammation. Surgery, typically sphincter-sparing (e.g., seton placement for drainage), may be used as an adjunct to medical therapy or as a secondary measure if medical treatment fails to achieve healing.
  \item \textbf{Complex or High Fistulas:} For complex or high trans-sphincteric, suprasphincteric, or extrasphincteric fistulas, a staged approach is common. Initial surgery may involve seton placement to drain the tract and allow inflammation to subside, followed by a definitive sphincter-preserving repair (e.g., advancement flap, LIFT procedure) at a later stage. In these cases, the initial surgery is a preparatory or salvage measure, with the definitive repair being the secondary, curative step.
  \item \textbf{Recurrent Fistulas:} For recurrent fistulas, surgery is a secondary intervention, often requiring more advanced imaging and complex techniques due to altered anatomy from previous procedures.
\end{itemize}

\section*{Surgical Technique \& Procedure}
Anal fistula surgery is performed under general or regional (spinal or caudal) anesthesia, typically with the patient in the prone jack-knife or lithotomy position to provide optimal exposure of the perianal region.

The general steps for anal fistula surgery include:
\begin{enumerate}
  \item \textbf{Examination Under Anesthesia (EUA):} A thorough digital rectal examination and anoscopy are performed to identify the external opening, palpate the tract, and locate the internal opening. Hydrogen peroxide, methylene blue dye, or milk may be injected into the external opening to help visualize the internal opening.
  \item \textbf{Tract Delineation:} A malleable probe is gently passed from the external opening through the tract to the internal opening. Care is taken to avoid creating false passages. The probe's position relative to the sphincter muscles is carefully assessed.
  \item \textbf{Fistulotomy (for low fistulas):} For intersphincteric or low trans-sphincteric fistulas involving less than 30-50\% of the external sphincter, the tissue overlying the probe is incised with electrocautery or a scalpel, laying the entire tract open. The edges of the wound are marsupialized to prevent premature skin closure.
  \item \textbf{Seton Placement (for high or complex fistulas):} For high trans-sphincteric, suprasphincteric, or extrasphincteric fistulas, or in cases of active inflammation (e.g., Crohn's), a seton (a non-absorbable suture or vessel loop) is passed through the tract and tied loosely. A draining seton promotes drainage and fibrosis, while a cutting seton is gradually tightened over weeks to slowly divide the sphincter muscle.
  \item \textbf{Endorectal Advancement Flap (for sphincter-preserving repair):} The internal opening is excised, and a flap of rectal mucosa and submucosa (and sometimes internal sphincter) is mobilized from above the internal opening. The internal opening is then closed primarily, and the flap is advanced distally and sutured over the repair, covering the internal opening. The external tract is curetted and drained.
  \item \textbf{Ligation of Intersphincteric Fistula Tract (LIFT Procedure):} An incision is made in the intersphincteric groove. The fistula tract is identified, ligated, and divided in the intersphincteric space. The internal opening is closed, and the external tract is curetted.
  \item \textbf{Fistula Plug or Fibrin Glue:} After curettage of the tract, a biodegradable plug (e.g., collagen matrix) is inserted into the tract, or fibrin glue is injected, aiming to promote tissue ingrowth and closure. These are typically used for less complex fistulas or as a secondary option.
  \item \textbf{Wound Management and Closure:} For fistulotomy, the wound is left open to heal by secondary intention. For sphincter-preserving techniques, the external wound may be left open or loosely closed, and drains may be placed if significant dead space exists.
\end{enumerate}

\section*{Preoperative Workup \& Preparation}
A comprehensive preoperative workup is essential to accurately characterize the fistula, assess the patient's overall health, and plan the most appropriate surgical approach.

\begin{itemize}
  \item \textbf{History and Physical Examination:} Detailed history regarding symptoms, previous abscesses, and medical comorbidities (e.g., inflammatory bowel disease, diabetes). A thorough perianal examination is crucial to identify external openings, palpate the tract, and assess sphincter tone.
  \item \textbf{Imaging Studies:} For complex, recurrent, or high fistulas, imaging is vital to delineate the tract's anatomy and its relationship to the sphincter complex, and to rule out associated abscesses.
    \begin{itemize}
      \item \textbf{Magnetic Resonance Imaging (MRI) of the pelvis:} Considered the gold standard for complex fistulas, providing excellent soft tissue detail.
      \item \textbf{Endoanal Ultrasound (EAUS):} Operator-dependent but useful for visualizing the sphincter muscles and the tract.
      \item \textbf{CT scan:} Less useful for fistula tracts but can identify large abscesses.
    \end{itemize}
  \item \textbf{Laboratory Tests:} Routine preoperative blood tests include a complete blood count (CBC), basic metabolic panel (BMP), coagulation profile (PT/INR, PTT), and blood type and screen.
  \item \textbf{Medical Optimization:} Patients with significant comorbidities (e.g., cardiac, pulmonary disease) may require medical clearance and optimization by their primary care physician or a specialist.
  \item \textbf{Bowel Preparation:} For most simple fistulas, a clear liquid diet for 24 hours prior to surgery and/or a cleansing enema on the morning of surgery is sufficient. For procedures involving an endorectal flap or extensive dissection, a more thorough mechanical bowel preparation may be required.
  \item \textbf{Medication Adjustments:} Anticoagulant and antiplatelet medications are typically held preoperatively according to institutional guidelines to minimize bleeding risk.
  \item \textbf{Antibiotic Prophylaxis:} Intravenous broad-spectrum antibiotics are usually administered within one hour prior to incision, especially for complex repairs or in patients with risk factors for infection (e.g., valvular heart disease, immunosuppression).
  \item \textbf{Informed Consent:} Detailed discussion with the patient regarding the planned procedure, alternative treatments, potential benefits, and risks, particularly the risk of fecal incontinence and recurrence.
\end{itemize}

\section*{Contraindications \& Precautions}
Careful consideration of contraindications and precautions is essential to optimize patient safety and surgical outcomes.

\textbf{Absolute Contraindications:}
\begin{itemize}
  \item Undrained acute perianal abscess: Any acute abscess must be drained and sepsis controlled before definitive fistula repair.
  \item Uncorrected severe coagulopathy or bleeding disorder: These must be corrected prior to surgery to minimize hemorrhagic risk.
  \item Severe, uncontrolled systemic sepsis: Patients must be medically stabilized before elective surgery.
\end{itemize}

\textbf{Relative Contraindications \& Precautions:}
\begin{itemize}
  \item \textbf{Active, severe Crohn's disease with extensive rectal inflammation:} Simple fistulotomy may not heal well and can lead to larger, non-healing wounds. Sphincter-sparing techniques or medical management are often preferred.
  \item \textbf{High trans-sphincteric, suprasphincteric, and extrasphincteric fistulas:} Simple fistulotomy is relatively contraindicated due to the high risk of significant sphincter division and subsequent fecal incontinence. Sphincter-preserving techniques or staged procedures with setons are indicated.
  \item \textbf{Pre-existing fecal incontinence or compromised sphincter function:} Any procedure involving further division of the sphincter muscles is relatively contraindicated. Careful assessment of continence preoperatively (e.g., manometry) is crucial, and sphincter-preserving techniques are mandatory.
  \item \textbf{Prior radiation to the pelvis or prior extensive anorectal surgery:} These can impair tissue healing and increase the risk of complications.
  \item \textbf{Immunosuppression (e.g., HIV/AIDS, organ transplant recipients):} Increased risk of infection and impaired wound healing.
  \item \textbf{Pregnancy:} Elective fistula surgery is generally deferred until after delivery, unless symptoms are severe and cannot be managed conservatively.
\end{itemize}

\section*{Complications \& Risks}
As with any surgical procedure, anal fistula surgery carries both general and procedure-specific risks.

\textbf{General Surgical Complications:}
\begin{itemize}
  \item \textbf{Bleeding:} Postoperative hemorrhage or hematoma formation, sometimes requiring re-intervention.
  \item \textbf{Infection:} Wound infection, cellulitis, or abscess formation, despite prophylactic antibiotics.
  \item \textbf{Pain:} Postoperative pain, which can be significant, requiring adequate analgesia.
  \item \textbf{Anesthesia risks:} Adverse reactions to anesthetic agents, respiratory or cardiac complications.
  \item \textbf{Deep vein thrombosis (DVT) and pulmonary embolism (PE):} Risk is generally low for perianal surgery but can occur.
  \item \textbf{Urinary retention:} Common after anorectal surgery, especially with spinal anesthesia.
\end{itemize}

\textbf{Procedure-Specific Complications \& Risks:}
\begin{itemize}
  \item \textbf{Fecal or flatus incontinence:} The most feared complication, directly related to the extent of sphincter muscle divided or damaged during surgery. This can range from minor flatus incontinence to severe fecal incontinence.
  \item \textbf{Fistula recurrence:} The fistula tract may fail to heal or may recur, sometimes with new tracts, requiring further surgery. Recurrence rates vary significantly by technique and fistula complexity.
  \item \textbf{Delayed wound healing:} Open wounds, particularly after fistulotomy, can take several weeks to months to heal completely.
  \item \textbf{Anal stricture:} Scarring after extensive surgery, especially fistulectomy or multiple procedures, can lead to narrowing of the anal canal, causing difficulty with defecation.
  \item \textbf{Persistent drainage:} Even after successful closure of the main tract, some patients may experience persistent serous drainage for a period.
  \item \textbf{Fecal impaction:} Can occur postoperatively due to pain, fear of defecation, or opioid use.
  \item \textbf{Damage to surrounding structures:} Rarely, in very complex or extrasphincteric cases, injury to the rectum, vagina, or urethra can occur.
  \item \textbf{Keyhole deformity:} After fistulotomy, the external opening may heal in a way that leaves a small indentation, which can trap stool and cause hygiene issues.
\end{itemize}

\section*{Postoperative Recovery Timeline}
The postoperative recovery timeline for anal fistula surgery varies significantly depending on the complexity of the fistula and the surgical technique employed.

\textbf{Immediate Postoperative Period (PACU and first 24-48 hours):}
Patients are monitored in the Post-Anesthesia Care Unit (PACU) for recovery from anesthesia, pain control, and vital sign stability. Most simple fistula surgeries are performed on an outpatient basis, with discharge home the same day. Pain is managed with oral analgesics, often a combination of NSAIDs and opioids. Patients are encouraged to ambulate early to prevent complications like DVT. The first bowel movement may be painful, and stool softeners are typically prescribed to prevent constipation and straining. Warm sitz baths (2-3 times daily) are highly recommended starting the day after surgery to promote hygiene, reduce pain and spasm, and encourage wound healing.

\textbf{First 1-2 Weeks Postoperatively:}
Patients are advised to maintain a high-fiber diet and adequate fluid intake to ensure soft stools. Activity restrictions usually include avoiding strenuous exercise, heavy lifting, and prolonged sitting for the first week or two. Wound care involves regular sitz baths and gentle cleaning. For open wounds (e.g., after fistulotomy), daily dressing changes or packing may be required, often performed by the patient or a visiting nurse. Pain typically subsides significantly within the first week, though discomfort with bowel movements may persist longer. Most patients can return to light work or daily activities within 1-2 weeks, depending on their comfort level and the nature of their work.

\textbf{Long-Term Recovery (Weeks to Months):}
Complete healing of an open fistulotomy wound can take anywhere from 4 weeks to 3 months, or even longer for very deep or complex tracts. Sphincter-preserving techniques generally have a shorter external wound healing time, but internal healing and integration of grafts/plugs take time. Follow-up appointments are crucial to monitor wound healing, assess for signs of infection or recurrence, and address any concerns regarding continence. Patients are advised to continue a high-fiber diet and avoid constipation long-term. Full recovery, including resolution of all symptoms and return to pre-surgical activity levels, can take several months.

\section*{Surgical Findings \& Pathology}
During anal fistula surgery, the surgeon meticulously documents several key findings. These include the precise location of the external opening(s), the course of the fistula tract in relation to the anal sphincter complex (e.g., intersphincteric, trans-sphincteric), the location of the internal opening (often at the dentate line), and the presence of any associated abscesses or secondary tracts. The surgeon also assesses the quality of the surrounding tissues and the integrity of the sphincter muscles.

Any excised tissue, such as the fistula tract itself or a portion of the internal opening, is sent for histopathological examination. The pathology report typically confirms the presence of a chronic inflammatory tract lined by granulation tissue or epithelial cells, consistent with a fistula-in-ano. It may also identify signs of chronic inflammation, fibrosis, and occasionally foreign body giant cells. In cases where Crohn's disease is suspected, the pathologist will look for characteristic granulomas, although their absence does not rule out the disease. The report helps confirm the diagnosis and excludes other pathologies, such as malignancy, which can mimic a fistula.

\section*{Expected Postoperative Course}
A normal and successful postoperative course following anal fistula surgery is characterized by a progressive reduction in pain and discomfort, followed by complete healing of the fistula tract. Patients should expect some pain, particularly with bowel movements, for the first few days to a week, which is manageable with prescribed oral analgesics. Drainage from the wound, initially serosanguinous, will gradually decrease and become clearer as healing progresses. For open wounds, the external opening will slowly close as granulation tissue fills the defect from the base upwards. There should be no signs of infection, such as increasing redness, swelling, fever, or purulent discharge. The patient should be able to resume a normal diet and gradually increase activity levels without significant limitations. Ultimately, the goal is complete cessation of drainage, closure of the external wound, and preservation of normal bowel and continence function, allowing the patient to return to their usual quality of life free from fistula-related symptoms.

\section*{Postoperative Care \& Follow-up}
Postoperative care is crucial for ensuring proper healing and detecting any complications or recurrence.
\begin{itemize}
  \item \textbf{Wound Care:} Patients are instructed on proper wound hygiene, including regular warm sitz baths (2-3 times daily) to keep the area clean, reduce pain, and promote healing. For open wounds, gentle packing or dressing changes may be required.
  \item \textbf{Pain Management:} Oral analgesics are prescribed, and patients are advised to take them as needed, especially before bowel movements.
  \item \textbf{Bowel Management:} A high-fiber diet, adequate fluid intake, and stool softeners are recommended to prevent constipation and straining, which can disrupt healing.
  \item \textbf{Activity Restrictions:} Patients are advised to avoid strenuous activities, heavy lifting, and prolonged sitting for several weeks, depending on the extent of surgery.
  \item \textbf{Follow-up Visits:} Scheduled appointments with the surgeon are essential to monitor wound healing, assess for signs of infection or recurrence, and evaluate continence. These visits typically occur at 1-2 weeks, 4-6 weeks, and then as needed until complete healing.
  \item \textbf{Warning Signs:} Patients are educated on warning signs that necessitate immediate contact with their surgeon, including increasing pain, fever, chills, worsening redness or swelling around the wound, purulent discharge, or significant bleeding.
\end{itemize}
In cases of complex fistulas or those associated with Crohn's disease, long-term follow-up and coordination with a gastroenterologist may be necessary.

\section*{Epidemiology}
Anal fistulas are a relatively common anorectal condition, with an estimated incidence ranging from 1 to 2 cases per 10,000 people per year in the general population. They predominantly affect adults, with the highest incidence observed in individuals between 30 and 50 years of age. There is a male predominance, with men being affected approximately 1.8 to 2 times more frequently than women. The vast majority (80-90\%) of anal fistulas are cryptoglandular in origin, arising from an infection of the anal glands that leads to an abscess and subsequent fistula formation. The most common types are intersphincteric and low trans-sphincteric fistulas. In patients with inflammatory bowel disease, particularly Crohn's disease, the lifetime risk of developing a perianal fistula is significantly higher, affecting up to 25-30\% of patients, and these fistulas tend to be more complex, recurrent, and challenging to treat. While mortality directly from anal fistula surgery is exceedingly rare, morbidity related to recurrence and incontinence remains a significant concern.

\section*{Outcomes \& Prognosis}
The outcomes and prognosis of anal fistula surgery are highly dependent on the fistula's complexity, its anatomical classification (e.g., Parks classification), the surgical technique employed, and the presence of underlying conditions like Crohn's disease. For simple, low intersphincteric or low trans-sphincteric fistulas treated with fistulotomy, the cure rate is excellent, often exceeding 90-95\%, with a low recurrence rate (typically less than 5-10\%) and a minimal risk of significant fecal incontinence.

However, for complex and high fistulas, the prognosis is more guarded. Sphincter-preserving techniques, such as endorectal advancement flaps, LIFT procedures, or fistula plugs, aim to reduce the risk of incontinence but may be associated with higher recurrence rates, often ranging from 10-30\%. Randomized controlled trials and large cohort studies consistently show a trade-off between continence preservation and cure rates for complex fistulas. Factors negatively impacting prognosis include high fistula tracts, multiple tracts, previous failed surgeries, and the presence of Crohn's disease. In patients with Crohn's-associated fistulas, long-term medical management is often crucial alongside surgical drainage to maintain healing and prevent recurrence, and the overall prognosis for complete, sustained healing is less predictable. Despite these challenges, modern surgical techniques, combined with careful patient selection and meticulous execution, offer good to excellent outcomes for the majority of patients, significantly improving their quality of life.

\section*{References}
\begin{enumerate}
  \item MedlinePlus. Anal fistula surgery. U.S. National Library of Medicine; 2024.
  \item Brunicardi FC, Andersen DK, Billiar TR, et al. Schwartz's Principles of Surgery. 11th ed. McGraw-Hill Education; 2019.
  \item American Society of Colon and Rectal Surgeons Clinical Practice Guidelines for the Management of Anorectal Abscess, Fistula-in-Ano, and Rectovaginal Fistula. Dis Colon Rectum. 2022;65(1):e1-e37.
  \item Parks AG, Gordon PH, Hardcastle PR. A classification of fistula-in-ano: with special reference to the intersphincteric and trans-sphincteric types. Br J Surg. 1976;63(1):1-12.
\end{enumerate}

\clearpage